%% kapa: exam=1 round=19 year=2008 period=2 form=A subject=economics subject_ko=경제원론 seq=20 number=20
%% source: Q-Net 감정평가사 기출문제 (cst00309) artlSeq=1001962 file=extracted/1008/1001962/2008년 제19회 감정평가사 1차시험 기출문제/2교시.hwp
%% conversion: kordoc 4.24.51 -> tools/gen1.py (auto); review=proofread-sonnet; flags=none
\begin{question}{20}
이윤극대화를 추구하는 어떤 기업의 생산함수가 $Y=A \cdot N^{ 1/2 } K^{ 1/2 }$(단, $A>0 , K>0 , N>0$)로 주어진다고 하자. 여기에서 $Y$ = 생산량, $N$ = 노동시간, $K$ = 자본투입량이며, $A$는 총요소생산성을 나타내는 상수이다. 이 생산함수와 부합하지 않는 것은? (단, 생산물시장과 생산요소시장은 완전경쟁적이라고 가정한다.)
\begin{choices}
  \item 총요소생산성이 커지면 노동시간 또는 자본투입량이 그대로 있더라도 생산량은 증가한다.
  \item 노동의 한계생산은 0보다 크나 노동시간이 증가할수록 노동의 한계생산은 점점 작아진다.
  \item 자본의 한계생산은 0보다 크지만 자본에 대해서는 한계생산체감의 법칙이 적용되지 않는다.
  \item 노동시간당 임금과 자본의 한 단위당 임대료가 동일하다면 노동시간과 자본투입량은 동일하게 해야 한다.
  \item 노동시간과 자본투입량을 각각 두 배로 늘리면 생산량도 두 배로 증가한다.
\end{choices}
\end{question}
\answer{3}
